\chapter{Conclusiones}

Para culminar este Trabajo de Fin de Grado, se revisará el cumplimiento de los objetivos planteados en la introducción y se presentarán varias líneas de trabajo futuro y propuestas de mejora, y una conclusión final.

\section{Cumplimiento de Objetivos}

En la siguiente tabla se resume el cumplimiento de los objetivos establecidos y la forma en que se han alcanzado.

\begin{table}[H]
    \centering
    \footnotesize
    \begin{tabularx}{\textwidth}{|p{0.2\textwidth}|p{0.65\textwidth}|c|}
    \hline
    \textbf{Objetivo} & \textbf{Cómo se ha conseguido} & \textbf{Estado} \\ \hline
    Sistema de autenticación & Registro con validación de email, inicio de sesión seguro con JWT, y recuperación de contraseña & \checkmark \\ \hline
    Bandeja de cartas enviadas y recibidas & Desarrollo de interfaces con filtrado por título, idioma, diario y visualización del estado de corrección & \checkmark \\ \hline
    Interfaz para crear, editar y enviar cartas & Diseño de editor intuitivo para especificar título, contenido, idioma, fecha y diario, guardar borradores y compartir con contactos & \checkmark \\ \hline
    Página de corrección de cartas & Página especializada con editor enriquecido, herramientas de adición y edición de correcciones y reenvío de cartas corregidas & \checkmark \\ \hline
    Sistema de solicitudes de amistad & Implementación de solicitudes de amistad, gestión de contactos, buscador de usuarios y algoritmo de recomendación de usuarios & \checkmark \\ \hline
    Página de perfil & Edición de datos personales, cambio de contraseña con verificación, eliminación de cuenta y carga segura de imagen & \checkmark \\ \hline
    Arquitectura limpia y escalable & Patrones de diseño reconocidos, separación de responsabilidades y optimización del frontend y del backend & \checkmark \\ \hline
    Metodología Scrum & Creación de historias de usuario, pruebas de validación, matriz de trazabilidad, y planificación por sprints, usando tableros de Trello & \checkmark \\ \hline
    Documentación completa & Redacción de todas las secciones de la presente memoria con contenido y forma adecuados  & \checkmark \\ \hline
    \end{tabularx}
    \caption{Cumplimiento de objetivos}
    \label{tab:cumplimiento_objetivos}
\end{table}

Todos los objetivos técnicos, de gestión y de documentación establecidos han sido alcanzados satisfactoriamente.

\section{Líneas de Trabajo Futuro y Propuestas de Mejora}

Tras la culminación de la fase de desarrollo actual, se han identificado diversas áreas que permitirían escalar la plataforma y mejorar la experiencia de usuario final. Estas propuestas se detallan a continuación:

\begin{itemize}
    \item \textbf{Migración a Almacenamiento en la Nube}: Actualmente, el sistema utiliza el almacenamiento local del servidor mediante Multer. Una mejora relevant consistiría en la integración de servicios como Amazon S3 o Cloudinary. Esto permitiría gestionar los recursos multimedia de forma más eficiente, mejorar su disponibilidad y escalabilidad, reducir la carga del servidor principal y estandarizar el despliegue de la aplicación.

    \item \textbf{Autenticación de Terceros}: Con el fin de reducir la fricción en el registro de nuevos usuarios, se plantea la implementación de Google Sign-In. Esto delegaría la seguridad de la autenticación a un proveedor de identidad consolidado, mejorando la confianza del usuario y la tasa de conversión de la plataforma.

    \item \textbf{Enriquecimiento de la Mensajería}: Evolucionar el sistema de cartas permitiendo adjuntar archivos de audio o imágenes dentro del cuerpo del mensaje. Esto potenciaría el aprendizaje lingüístico al permitir que los usuarios practiquen no solo la expresión escrita, sino también la comprensión auditiva y la pronunciación.

    \item \textbf{Módulo de Comunidad}: Adición de una página de comunidad donde los usuarios puedan crear grupos para compartir cartas, temas sobre los que escribir, foros, canciones, etcétera. Esto podría mejorar la experiencia del usuario al tener acceso a más funcionalidades y a una conexión más grupal con otros usuarios; aunque antes de su implementación se debería realizar un estudio para determinar si existiría un interés real en esta nueva página por parte de los usuarios de la app.
\end{itemize}


\section{Conclusión final}

La realización de este Trabajo de Fin de Grado ha representado un hito fundamental en mi formación académica. El desarrollo de una plataforma full-stack para el intercambio lingüístico me ha permitido aplicar de manera transversal y práctica conocimientos teóricos adquiridos durante el grado, y profundizar en tecnologías de interés personal y profesional.

En el reto de construir un ecosistema de software completo, robusto y funcional desde sus cimientos, la consolidación de habilidades técnicas ha ido de la mano con el aprendizaje autónomo. La investigación del marco legislativo pertinente y la implementación de medidas de seguridad avanzadas para la protección contra ciberataques me ha llevado a un conocimiento más profundo y práctico los estándares actuales de la industria, comprendiendo que el desarrollo de software no solo consiste en crear funcionalidad, sino en garantizar la integridad y privacidad de los datos del usuario.

Asimismo, este trabajo ha sido un ejercicio de resiliencia y resolución de problemas. Encontrarme con desafíos técnicos, como la gestión de archivos estáticos con Multer, la autenticación mediante tokens JWT o la sincronización del estado global en Next.js, me ha permitido desarrollar una mentalidad analítica para buscar soluciones eficientes y documentadas. Esta capacidad de "aprender a aprender" es, sin duda, la competencia más valiosa que me llevo para mi futura etapa profesional.

Desde el punto de vista de la gestión, la aplicación de la metodología Scrum y el uso de herramientas como Trello han sido determinantes. La necesidad de reajustar los sprints y organizar las tareas de forma estructurada, funcional y centrada en el aporte de valor me ha proporcionado una visión realista de la ingeniería de software: un equilibrio constante entre requisitos, tiempo y calidad. Se optó por el estudio y uso de esta metodología sobretodo para el aprendizaje y acercamiento a formas de trabajo reales de la industria. Sin embargo, se ha observado que no es el mejor enfoque para proyectos desarrollados principalmente por una sola persona, dado que está más orientado a entornos colaborativos con múltiples roles y equipos.

En definitiva, este proyecto no solo ha validado mi capacidad técnica como futura graduada, sino que ha forjado mi identidad como ingeniera informática, dotándome de la confianza y las herramientas necesarias para abordar los retos tecnológicos reales de la industria.